\documentclass[12pt,a4paper]{article}

% 基础包
\usepackage{ctex}        % 中文支持
\usepackage{geometry}     % 页面设置
\usepackage{fancyhdr}     % 页眉页脚
\usepackage{lastpage}     % 总页数
\usepackage{indentfirst}  % 首行缩进

% 字体设置 - 使用Times New Roman字体（含数字）
\usepackage{mathptmx}     % 数学字体为Times New Roman风格
\usepackage[T1]{fontenc}  % 字体编码

% 数学公式
\usepackage{amsmath,amssymb,amsfonts}
\usepackage{mathtools}
\usepackage{nccmath}
\everymath{\displaystyle}  % 全文行内公式自动启用 displaystyle
\usepackage{txfonts}

% 表格
\usepackage{array}
\usepackage{booktabs}
\usepackage{multirow}
\usepackage{colortbl}

% 图片和颜色
\usepackage{graphicx}
\usepackage{xcolor}
\usepackage{caption}
\usepackage{subcaption}

% 列表
\usepackage{enumitem}
\usepackage{tasks}        % 用于排版选择题选项

% 页面设置
\geometry{
	left=1.6cm,
	right=2.04cm,
	top=1.6cm,
	bottom=2.54cm,
	headheight=1.2cm,
	footskip=0.8cm,
	nohead                 % 去掉页眉
}

% 行距和段落 - 1.5倍行距
\linespread{1.5}
\setlength{\parindent}{2em}  % 首行缩进2字符
\setlength{\parskip}{0.5ex}

% 页眉页脚设置 - 只有页脚，没有页眉，去掉横线
\pagestyle{fancy}
\fancyhf{}                    % 清空所有页眉页脚
\fancyfoot[C]{\makebox[0pt][c]{数学试题~第\thepage 页（共\pageref{LastPage}页）}}  % 使用\makebox确保居中对齐
\renewcommand{\headrulewidth}{0pt}  % 去掉页眉横线
\renewcommand{\footrulewidth}{0pt}  % 去掉页脚横线

% 选择题选项环境配置
% 4个选项一行（用于单选题）
\NewTasksEnvironment[
label=\Alph*.,
label-width=1.5em,
item-indent=2em,
column-sep=2em,
before-skip=0.8em,
after-skip=0.8em
]{choicesFour}[\choice](4)

% 2个选项一行（用于多选题和部分单选题）
\NewTasksEnvironment[
label=\Alph*.,
label-width=1.5em,
item-indent=2em,
column-sep=3em,
before-skip=0.8em,
after-skip=0.8em
]{choicesTwo}[\choice](2)

% 1个选项一行
\NewTasksEnvironment[
label=\Alph*.,
label-width=1.5em,
item-indent=2em,
before-skip=0.8em,
after-skip=0.8em
]{choicesOne}[\choice](1)

\begin{document}
	
	% 左上角绝密标识 - 使用$\bigstar$命令确保显示
	\noindent {\large\bfseries 绝密$\bigstar$启用前}
	
	\vspace*{0.5cm}
	
	% 封面部分
	\begin{center}
		{\Large 2026年普通高等学校招生全国统一考试}\\[0.8cm]
		{\LARGE\bfseries 数~~~~学}
	\end{center}
	
	\noindent 本试卷共4页，19小题，满分150分，考试时间120分钟.
	
	% 注意事项 - 严格按照您提供的格式
	\noindent\hangindent=73pt {\bfseries 注意事项：}1. 答题前，请务必将自己的姓名、准考证号用黑色字迹签字笔或钢笔分别填写在试卷和答题卡上.  将条形码横贴在答题卡右上角的“条形码粘贴处”. 
	
	\noindent\hangindent=73pt\hspace{61pt}2. 作答选择题时，选出每小题答案后，用2B铅笔把答题卡上对应题目选项的答案信息点涂黑，在试卷上作答无效；如需改动，用橡皮擦干净后，再选涂其他答案. 
	
	\noindent\hangindent=73pt\hspace{61pt}3. 非选择题必须用黑色字迹钢笔或签字笔作答，答案必须写在答题卡各题目指定区域的相应位置上；如需改动，先划掉原来的答案，然后再写上新的答案. 
	
	\noindent\hangindent=73pt\hspace{61pt}4. 考生须保持答题卡的整洁. 考试结束后，将试卷和答题卡一并交回. 
	
	
	% 一、选择题
	\noindent\hangindent=2em {\bfseries 一、选择题:本题共8小题，每小题5分，共40分.每小题给出的四个选项中，只有一项是符合题目要求的.}
	
	\begin{enumerate}[leftmargin=*,label=\arabic*.,itemsep=0.8ex]
		\item 样本数据$6, 8, 4, 5, 12$的中位数为
		
		\begin{choicesFour}
			\choice $5$ \choice $6$ \choice $8$ \choice $9$
		\end{choicesFour}
		
		\item 已知平面向量$\boldsymbol{a},\boldsymbol{b}$不共线，且$2\boldsymbol{a} + y\boldsymbol{b} = x\boldsymbol{a} - 3\boldsymbol{b}$，则
		
		\begin{choicesFour}
			\choice $x=2,\ y=-3$ \choice $x=-2,\ y=3$ \choice $x=2,\ y=3$ \choice $x=-2,\ y=-3$
		\end{choicesFour}
		
		\item 已知集合$A=\left\{\sin\frac{7\pi}{3},\cos\frac{5\pi}{4},\tan\frac{5\pi}{4}\right\},\ B=\left[-\frac{\sqrt{3}}{2},-\frac{1}{2}\right]$，则$A\cap B=$
		
		\begin{choicesFour}
			\choice $\left(-\frac{\sqrt{3}}{2},\frac{1}{2}\right)$ \choice $\left(-\frac{\sqrt{3}}{2},1\right)$ \choice $\left[-\frac{1}{2},1\right)$ \choice $\left(-\frac{\sqrt{3}}{2},-\frac{1}{2}\right)$
		\end{choicesFour}
		
		\item 曲线$y=5x+8\ln x$在点$(1,5)$处的切线方程为
		
		\begin{choicesFour}
			\choice $y=3x+2$ \choice $y=5x$ \choice $y=8x-3$ \choice $y=13x-8$
		\end{choicesFour}
		
		\item 已知抛物线$C_1:y^2=2p_1x\ (p_1>0)$和$C_2:x^2=2p_2y\ (p_2>0)$均经过点$(4,8)$，则$C_1$的焦点与$C_2$的焦点之间的距离为
		
		\begin{choicesFour}
			\choice $12$ \choice $4\sqrt{5}$ \choice $6$ \choice $\frac{\sqrt{65}}{2}$
		\end{choicesFour}
		
		\item 已知函数$f(x)=\frac{x+2}{e^x+a}$的最大值为$1$，则$a=$
		
		\begin{choicesFour}
			\choice $\frac{1}{2}$ \choice $1$ \choice $\frac{3}{2}$ \choice $2$
		\end{choicesFour}
		
		\item 一百零八塔位于宁夏回族自治区青铜峡市，其独特的建筑格局和深远的历史文化闻名遐迩，该塔群共有$108$座塔，依山势自上而下排成$12$行，将第$i$行中塔的座数记为$a_i\ (i=1,2,\dots,12)$，其中$a_1=1,\ a_3=a_2=3,\ a_8=a_5=5$，且$a_4,a_7,\dots,a_{12}$是一个首项为$7$，公差为$2$的等差数列。将$a_1,a_4,a_7,\dots,a_{12}$分为$6$组，每组$2$个数，使得每组的$2$个数之和可构成一个项数为$6$且公差为$d\ (d>0)$的等差数列，则$d=$
		
		\begin{choicesFour}
			\choice $2$ \choice $4$ \choice $6$ \choice $8$
		\end{choicesFour}
		
		\item 设$U=\{(x_1,x_2,x_3)\mid x_i\in\{-2,-1,1,2\},i=1,2,3\}$为空间中$64$个点构成的集合，点$P(1,1,1)$，记样本空间$\Omega=U\setminus\{P\}$。从$\Omega$中随机取一个点，定义随机变量$X$如下：对$\Omega$中的每个点$A(x_1,x_2,x_3)$，令$X(A)=x_1+x_2+x_3$，则$X$的数学期望值为
		
		\begin{choicesFour}
			\choice $-\frac{1}{21}$ \choice $-\frac{1}{63}$ \choice $0$ \choice $\frac{1}{7}$
		\end{choicesFour}
	\end{enumerate}
	
	\vspace{0.3cm}
	
	% 二、选择题
	\noindent\hangindent=2em {\bfseries 二、选择题:本题共3小题，每小题6分，共18分.每小题给出的选项中，有多项符合题目要求. 全部选对的得6分，部分选对的得部分分，有选错的得0分.}
	
	\begin{enumerate}[leftmargin=*,label=\arabic*.,resume,itemsep=0.8ex]
		\item 设$z=3+2\mathrm{i}$，则
		
		\begin{choicesTwo}
			\choice $\overline{z}=3-2\mathrm{i}$ \choice $|z|=5$ \choice $z^2=5+12\mathrm{i}$ \choice $\frac{z+3}{z-1}\in\mathbb{R}$
		\end{choicesTwo}
		
		\item 在空间中，$A,B$为两个定点，动点$C$到直线$AB$的距离为$2$，动点$D$到直线$AB$的距离为$1$。若二面角$C-AB-D$为$60^\circ$，则
		
		\begin{choicesOne}
			\choice $\angle CAD\ge60^\circ$ \choice $CD\ge\sqrt{3}$ \choice 当$AB\perp CD$时，$CD\perp$平面$ABD$ \choice 当$AB\perp$平面$ACD$时，$AC\perp AD$
		\end{choicesOne}
		
		\item 已知圆$C_1:(x+1)^2+y^2=1$，圆$C_2:(x-1)^2+y^2=1$，圆$C_3:x^2+(y-\sqrt{3})^2=1$，直线$l:y=kx+b$与$C_1,C_2,C_3$均有两个交点，记$l$被$C_1,C_2,C_3$截得的弦长分别为$s_1,s_2,s_3$，则
		
		\begin{choicesOne}
			\choice $k$可以取任意实数 \choice 满足$s_1=s_2=s_3$的直线$l$共有$3$条 \choice 满足$s_1+s_2+s_3=3$的直线$l$多于$3$条 \choice 当$b=0$时，$s_1+s_2+s_3$的最大值为$\frac{2\sqrt{21}}{3}$
		\end{choicesOne}
	\end{enumerate}
	
	\vspace{0.3cm}
	
	% 三、填空题
	\noindent {\bfseries 三、填空题:本题共3小题，每小题5分，共15分.}
	
	\begin{enumerate}[leftmargin=*,label=\arabic*.,resume,itemsep=0.8ex]
		\item 双曲线$5x^2-6y^2=1$的离心率为\underline{\qquad\qquad}. 
		
		\item 已知$f(x)=2\sin(\omega x+\theta)\ (\omega\in\mathbb{Z},\ 0\le\theta<2\pi)$是偶函数，$f(x)$在区间$\left(0,\frac{\pi}{2}\right)$单调递增，则$\theta=$\underline{\qquad\qquad}，$f\left(\frac{2\pi}{3}\right)=$\underline{\qquad\qquad}. 
		
		\item 设实数$q$满足：存在数列$\{a_n\}$，使得对于任意$n\in\mathbb{N}^*$，均有$a_1+a_2+\dots+a_n=n^2+n$，且$\{a_n\}$中有某连续$9$项$a_k,a_{k+1},\dots,a_{k+8}$是公比为$q$的等比数列，则$q$的最大值为\underline{\qquad\qquad}. 
	\end{enumerate}
	
	\vspace{0.3cm}
	
	% 四、解答题
	\noindent {\bfseries 四、解答题:本题共5小题，共77分.解答应写出文字说明、证明过程或演算步骤.}
	
	\begin{enumerate}[leftmargin=*,label=\arabic*.,resume,itemsep=1.2ex]
		\item (13分)\\
		如图，在直三棱柱$ABC-A_1B_1C_1$中，$\angle ACB=90^\circ$，$AC=BC$，$D,E$分别为$AB,AC_1$的中点.
		
		(1) 证明：$DE\varparallel$平面$BCC_1B_1$；
		
		(2) 设$CC_1=2$，直线$DE$与平面$ACC_1A_1$所成的角为$45^\circ$，求直线$DE$到平面$BCC_1B_1$的距离. 
		
		\item (15分)\\
		已知在$\triangle ABC$中，$AB=3$，$BC=2\sqrt{3}$，$\cos B=\frac{\sqrt{3}}{3}$.
		
		(1) 求$\cos A$；
		
		(2) 设$D,E$两点满足：$D$在$BA$的延长线上，$DE\varparallel BC$，$AE\perp AC$. 若$DE=\sqrt{6}$，求$CE$. 
		
		\item (15分)\\
		设整数$N\ge2$. 某同学用一个球进行投篮练习，至多投篮$N$次，当且仅当投中1次时或$N$次均未投中时，停止练习. 设该同学每次投中的概率为$p\ (0<p<1)$，各次投中与否相互独立. 记$X$为停止练习时该同学的投篮次数.
		
		(1) 当$\displaystyle N=4,\ p=\frac{1}{3}$时，求$X$的分布列；
		
		(2) 设$k,m$均为自然数.
		
		\ \ (i) 当$k\le N-1$时，求$P(X>k)$；
		
		\ \ (ii) 当$k+m\le N-1$时，证明：$P(X>k+m\mid X>k)=P(X>m)$. 
		
		\item (17分)\\
		已知椭圆$\displaystyle C:\frac{x^2}{a^2}+\frac{y^2}{b^2}=1\ (a>b>0)$的左焦点为$F(-1,0)$，离心率为$\displaystyle \frac{1}{2}$.
		
		(1) 求$C$的方程；
		
		(2) 设$O$为坐标原点，过$F$且斜率大于$0$的动直线$l$与$C$交于$P,Q$两点，其中$Q$在第三象限，直线$PO$与$C$的另一个交点为$R$.
		
		\ \ (i) 若$\triangle PQR$的面积是$\triangle PFO$的面积的$3$倍，求$l$的方程；
		
		\ \ (ii) 求$\tan\angle PQR$的最小值. 
		
		\item (17分)\\
		已知函数$f(x)$的定义域为$\mathbb{R}$，且当$x<0$时，$f(x)=2^x$. 对任意$x_0\in\mathbb{R}$，定义集合
		\[
		D(x_0)=\{d\in\mathbb{R}\mid f(x_0+d)>f(x_0)\}.
		\]
		
		(1) 若当$x\ge0$时，$f(x)=1-x$，求$D(-1)$；
		
		(2) 若$f(x)$是奇函数，$f(x_1)\le f(x_2)$，且$x_1x_2\ne0$，证明：$D(x_1)\subseteq D(x_2)$；
		
		(3) 设$f(x)$满足：① 若$f(x_1)\le f(x_2)$，则$D(x_2)\subseteq D(x_1)$；② 当$0<x<1$时，$f(x)<f(0)$.
		
		\ \ (i) 证明：$f(0)\ge1$；
		
		\ \ (ii) 证明：$f(x)$在区间$(0,+\infty)$单调递增. 
	\end{enumerate}
	
\end{document}